% `tex.web` §1090: `\-` at the head of a paragraph starts one.
%
% §1116 registers `\discretionary` and `\-` as one command with two `chr`
% values, and §1090's list of commands that begin a paragraph names that
% command once:
%
%     vmode+…,vmode+discretionary,…:
%       begin back_input; new_graf(true); end;
%
% We had only the braced form in the list, so trip.tex's `\--\--\char-0-A\-`
% left the first `\-` in vertical mode. It appended its discretionary to the
% enclosing vertical list, where no line breaker will ever look at it, and the
% paragraph began at the character after it instead — one node in the wrong
% list and an indent box in the wrong place.
%
% §1117 is what the node is:
%
%     tail_append(new_disc);
%     if cur_chr=1 then
%       begin c:=hyphen_char[cur_font];
%       if c>=0 then if c<256 then pre_break(tail):=new_character(cur_font,c);
%       end
%
% so `\-` is a discretionary whose pre-break text is the current font's
% \hyphenchar, and with \hyphenchar negative it is an empty one — still a
% legal place to break, just one that shows nothing.
\catcode`\{=1 \catcode`\}=2
\tracingonline=1 \showboxbreadth=99 \showboxdepth=99
\hsize=200pt \parindent=0pt \vsize=200pt
\font\r=cmr10 \r \hyphenchar\r=`\-

% `\-` first: it is what starts the paragraph, so the indent box is above it.
\setbox0\vbox{\-AB\par}
\showbox0

% ...and in the middle and at the end, where it never had to start anything.
\setbox1\vbox{A\-B\par}
\showbox1
\setbox2\vbox{A\-\par}
\showbox2

% A negative \hyphenchar gives an empty discretionary, not no discretionary.
\hyphenchar\r=-1
\setbox3\vbox{\-AB\par}
\showbox3

% In restricted horizontal mode there is no paragraph to start and the node is
% appended where it stands.
\hyphenchar\r=`\-
\setbox4\hbox{A\-B}
\showbox4
\end
